
\chapter{Derived categories}\label{VIII}
\blfootnote{by Jean-Louis Verdier}

\section{Triangulated categories: definitions and examples}\label{VIII:1}

\subsection{Definition of triangulated categories}\label{VIII:1-1}

\addtocounter{subsubsection}{-1}
\subsubsection{}\label{VIII:1-1-0}

Let $\cA$ be an additive category and let $T$ be an additive automorphism of
$\cA$. For two objects $X$ and $Y$ of $\cA$, we set
\[
  \hom^i(X,Y) = \hom_\cA(X,T^i(Y)) \qquad i\in \dZ
\]
Let $X$, $Y$, and $Z$ be three objects of $\cA$, let
$\alpha\in \hom^i(X,Y)$, and let $\beta\in \hom^j(Y,Z)$. We define the
composite $\beta\circ\alpha$ belonging to $\hom^{i+j}(X,Z)$ by
\[
  \beta\circ \alpha = T^i(\beta)\circ\alpha \text{.}
\]
One checks that this gives a new category whose groups of morphisms between
objects are graded. By an abuse of language, this new category will be called
the category $\cA$ graded by the translation functor $T$.

Let $\cA$ be a category graded by a translation functor $T$. Whenever a
morphism is named or denoted without its degree being specified, it will
always mean a morphism of degree zero.

Let $\cA$ and $\cA'$ be two additive categories graded by the functors $T$
and $T'$, and let $F$ be an additive functor from $\cA$ to $\cA'$. We say
that $F$ \emph{is graded} if an isomorphism of
functors\footnote{Text amended in 1976: the original text required an
equality $F\circ T=T'\circ F$. The case of functors in several variables is
discussed in \cite[XVIII 0.2]{sga4}. It raises a problem of signs.}
has been specified:
\[
  F\circ T = T'\circ F
\]

The definition of morphisms of graded functors is given in
\ref{VIII:4-1-1}.

Let $\cA$ be an additive category graded by the functor $T$. A triangle is
a collection of three objects $X$, $Y$, $Z$ and three morphisms
$u:X\to Y$, $v:Y\to Z$, $w:Z\to T(X)$ ($\deg w=1$). We use either the
notation $(X,Y,Z,u,v,w)$ or the following diagram for a triangle:
\[\xymatrix{
  & Z \ar[dl]_-w \\
  X \ar[rr]^-u
    & & Y \ar[lu]_-v
}\qquad \deg(w)=1
\]
Let $(X,Y,Z,u,v,w)$ and $(X',Y',Z',u',v',w')$ be two triangles of $\cA$.
A morphism from the first triangle to the second is a collection of three
morphisms $f:X\to X'$, $g:Y\to Y'$, $h:Z\to Z'$ such that the following
diagram is commutative:
\[\xymatrix{
  X \ar[r]^-u \ar[d]^-f
    & Y \ar[r]^-v \ar[d]^-g
    & Z \ar[r]^-w \ar[d]^-h
    & T(X) \ar[d]^-{T(f)} \\
  X' \ar[r]^-{u'}
    & Y' \ar[r]^-{v'}
    & Z' \ar[r]^-{w'}
    & T(X')
}\]

\subsubsection{}\label{VIII:1-1-1}

A \emph{triangulated category} is an additive category graded by a
translation functor and equipped with a family of triangles, called the
family of distinguished triangles. This family must also satisfy the
following axioms.

\paragraph{TR1}
\phantomsection\hypertarget{VIII:TR1}{}

Every triangle isomorphic to a distinguished triangle is distinguished.
Every morphism $u:X\to Y$ is contained in a distinguished triangle
$(X,Y,Z,u,v,w)$. The triangle
$(X,X,0,\operatorname{id}_X,0,0)$ is distinguished.

\paragraph{TR2}
\phantomsection\hypertarget{VIII:TR2}{}

A triangle $(X,Y,Z,u,v,w)$ is distinguished if and only if
$(Y,Z,T(X),v,w,-T(u))$ is distinguished.

\paragraph{TR3}
\phantomsection\hypertarget{VIII:TR3}{}

If $(X,Y,Z,u,v,w)$ and $(X',Y',Z',u',v',w')$ are distinguished, then for
every morphism $(f,g):u\to u'$, there exists a morphism $h:Z\to Z'$ such
that $(f,g,h)$ is a morphism of triangles.

\paragraph{TR4}
\phantomsection\hypertarget{VIII:TR4}{}

\[\xymatrix{
  & Z' \ar[dl]|{\deg=1} \\
  X \ar[rr]^-u
    & & Y \ar[ul]_-i
}\quad
\xymatrix{
  & X' \ar[dl]|{\deg(j)=1} \\
  Y \ar[rr]^-v
    & & Z \ar[ul]
}\quad
\xymatrix{
  & Y' \ar[dl]|{\deg=1} \\
  X \ar[rr]^-w
    & & Z \ar[ul]
}\]
If these are three distinguished triangles such that $w=v\circ u$, then
there exist two morphisms $f:Z'\to Y'$, $g:Y' \to X'$ such that
\begin{enumerate}
  \item $(\operatorname{id}_X, v, f)$ is a morphism of triangles;
  \item $(u,\operatorname{id}_Z,g)$ is a morphism of triangles;
  \item $(Z',Y',X',f,g,T(i) j)$ is a distinguished triangle.
\end{enumerate}

If $\cC$ and $\cC'$ are two triangulated categories, an \emph{exact
functor} from $\cC$ to $\cC'$ is an additive graded functor that takes
distinguished triangles to distinguished triangles. Morphisms of exact
functors are morphisms of graded functors.

The following property follows immediately from axioms
\hyperlink{VIII:TR1}{TR1}, \hyperlink{VIII:TR2}{TR2}, and
\hyperlink{VIII:TR3}{TR3}.

\begin{proposition}\label{VIII:1-1-2}
Let $(X,Y,Z,u,v,w)$ be a distinguished triangle and let $M$ be an object
of a triangulated category. Then the following sequences are exact:
\[\xymatrix{
  \cdots \ar[r]
    & \hom^i(M,X) \ar[r]^-{\hom(M,u)}
    & \hom^i(M,Y) \ar[r]^-{\hom(M,v)}
    & \hom^i(M,Z) \ar[r]^-{\hom(M,w)}
    & \hom^{i+1}(M,X) \ar[r]
    & \cdots
}\]
\[\xymatrix{
  \cdots \ar[r]
    & \hom^i(Z,M) \ar[r]^-{\hom(v,M)}
    & \hom^i(Y,M) \ar[r]^-{\hom(u,M)}
    & \hom^i(X,M) \ar[r]^-{\hom(w,M)}
    & \hom^{i+1}(Z,M) \ar[r]
    & \cdots
}\]
\end{proposition}

From this proposition one obtains statements of the following kind:
\begin{itemize}
  \item If $(f,g,h)$ is a morphism of distinguished triangles and $f,g$
    are isomorphisms, then $h$ is an isomorphism.
  \item The distinguished triangles constructed on a morphism
    (\hyperlink{VIII:TR1}{TR1}) are all isomorphic (\emph{but the
    isomorphisms are not, in general, uniquely determined}).
  \item When one of the morphisms in a distinguished triangle has a kernel
    or a cokernel, the latter splits; that is, the triangle is of the form
    \[\xymatrix{
      & Z+T(X) \ar[dl] \\
      X+Y \ar[rr]
        & & Y+Z \ar[ul]
    }\]
  \item The sum of two distinguished triangles is distinguished. (One uses
    axiom \hyperlink{VIII:TR4}{TR4}.)
\end{itemize}

\subsection{Examples}\label{VIII:1-2}

\subsubsection{}\label{VIII:1-2-1}

We first fix some notation. Let $\cA$ be an additive category.
The following category will be denoted by $\eC(\cA)$:
\begin{itemize}
  \item The objects of $\eC(\cA)$ are the complexes of $\cA$, with no
    restriction on degree and with differential of degree $+1$.
  \item The morphisms of $\eC(\cA)$ are the degree-preserving morphisms of
    complexes (those commuting with the differential).
\end{itemize}

Let $T$ be the following functor. For every object $X^\bullet$ of
$\eC(\cA)$,
\begin{align*}
  T(X^\bullet)_i &= (X^\bullet)_{i+1} \\
  d_i(T(X^\bullet)) &= -d_{i+1}(X^\bullet)  \text{.}
\end{align*}
On morphisms the functor $T$ acts as follows: for every morphism
$f:X^\bullet \to Y^\bullet$, the morphism
$T(f):T(X^\bullet) \to T(Y^\bullet)$ is defined by
\[
  T(f)_i = f_{i+1} \text{.}
\]
It is immediate that $T$ is an additive automorphism of $\eC(\cA)$.
The morphisms of degree $n$ defined using this translation functor are
then the morphisms that raise degree by $n$ and commute or anticommute
with the differential according to the parity of $n$. Thus one recovers
the usual definition.

Let $\sC$ denote the following family of triangles. A triangle
$(X^\bullet,Y^\bullet,Z^\bullet,u,v,w)$ belongs to the family
if\footnote{Text amended in 1976 to ensure the validity of
\ref{VIII:1-2-4}.}
\begin{itemize}
  \item $Z^\bullet$ is the simple complex associated with the double
    complex
    \[\xymatrix{
      \cdots \ar[r]
        & 0 \ar[r]
        & X^\bullet \ar[r]^-u
        & Y^\bullet \ar[r]
        & 0 \ar[r]
        & \cdots
    }\]
    in which the objects of $Y^\bullet$ have first degree zero and the
    objects of $X^\bullet$ have first degree $-1$.
  \item $v$ is the image, under the functor: double complex $\mapsto$
    simple complex, of the following morphism of double complexes:
    \[\xymatrix{
      \cdots \ar[r]
        & 0 \ar[r] \ar[d]
        & 0^\bullet \ar[r] \ar[d]
        & Y^\bullet \ar[r] \ar[d]
        & 0 \ar[r] \ar[d]
        & \cdots \\
      \cdots \ar[r]
        & 0 \ar[r]
        & X^\bullet \ar[r]
        & Y^\bullet \ar[r]
        & 0 \ar[r]
        & \cdots
    }\]
  \item $w$ is the negative of the image, under the same functor, of the
    following morphism of double complexes:
    \[\xymatrix{
      \cdots \ar[r]
        & 0 \ar[r]
        & X^\bullet \ar[r]
        & 0 \ar[r]
        & 0 \ar[r]
        & \cdots \\
      \cdots  \ar[r]
        & 0 \ar[r] \ar[u]
        & X^\bullet \ar[r] \ar[u]
        & Y^\bullet \ar[r] \ar[u]
        & 0 \ar[r] \ar[u]
        & \cdots
    }\]
\end{itemize}

\subsubsection{}\label{VIII:1-2-2}

The set of null-homotopic morphisms is a two-sided ideal (if $f$ and $g$
are composable and either $f$ or $g$ is null-homotopic, then the composite
is null-homotopic; if $f$ and $g$ are null-homotopic, then
$f\oplus g$ is null-homotopic). For every pair of objects $X^\bullet$,
$Y^\bullet$, let $\operatorname{Htp}(X^\bullet,Y^\bullet)$ denote the
subgroup of $\hom_{\eC(\cA)}(X^\bullet,Y^\bullet)$ consisting of the
null-homotopic morphisms. The objects of $\eK(\cA)$ are then the objects
of $\eC(\cA)$, and its group of morphisms with source $X^\bullet$ and
target $Y^\bullet$ is
\[
  \hom_{\eC(\cA)}(X^\bullet,Y^\bullet) / \operatorname{Htp}(X^\bullet,Y^\bullet) \text{.}
\]
The composition law on the morphisms of $\eK(\cA)$ is defined by passing
the composition law on the morphisms of $\eC(\cA)$ to the quotient.
The category $\eK(\cA)$ is additive. The translation functor $T$ on
$\eC(\cA)$ passes to the quotient (the translate of a null-homotopic
morphism is null-homotopic). The resulting functor is an additive
automorphism of $\eK(\cA)$, again denoted by $T$. By an abuse of notation,
$\eK(\cA)$ will henceforth denote the category $\eK(\cA)$ graded by the
functor $T$.

A distinguished triangle in $\eK(\cA)$ is any triangle \emph{isomorphic}
to a triangle arising from $\sC$. The family of distinguished triangles
satisfies axioms \hyperlink{VIII:TR1}{TR1},
\hyperlink{VIII:TR2}{TR2}, \hyperlink{VIII:TR3}{TR3}, and
\hyperlink{VIII:TR4}{TR4}. By an abuse of notation, $\eK(\cA)$ will denote
the category $\eK(\cA)$ triangulated by the family of triangles defined
above.

\subsubsection{}\label{VIII:1-2-3}

There is also a whole arsenal of subcategories:
\begin{itemize}
  \item A complex is said to be bounded below if all but finitely many of
    its objects of negative degree are zero. Bounded-above and bounded
    complexes are defined similarly. Let $\eK^+(\cA)$, $\eK^-(\cA)$,
    and $\eK^b(\cA)$ denote the full subcategories of $\eK(\cA)$ generated
    by these families of objects. These subcategories are
    “\emph{triangulated subcategories}” of $\eK(\cA)$: every distinguished
    triangle two of whose objects belong to the subcategory is isomorphic
    to a triangle all three of whose objects belong to the subcategory.
  \item Suppose that $\cA$ is an abelian category. A complex is said to be
    cohomologically bounded below if all but finitely many of its
    cohomology objects of negative degree are zero. Cohomologically
    bounded-above complexes, cohomologically bounded complexes, and
    acyclic complexes are defined similarly. Let
    $\eK^{\infty,+}(\cA)$, $\eK^{\infty,-}(\cA)$,
    $\eK^{\infty,b}(\cA)$, and $\eK^{\infty,\varnothing}(\cA)$ denote the
    full subcategories generated by these families of objects. They are
    triangulated subcategories.

    We shall also use the triangulated subcategories $\eK^{+,b}(\cA)$,
    $\eK^{+,\varnothing}(\cA)$, \ldots etc \ldots (the first superscript
    gives information about the objects of the complexes, and the second
    about the cohomology objects).
  \item Without assuming that $\cA$ is abelian, let $O$ be a set of objects
    of $\cA$ stable under isomorphisms and direct sums. Let $O(\cA)$ denote
    the full subcategory of $\eC(\cA)$ generated by the complexes whose
    objects in every degree belong to $O$. Let $\underline O(\cA)$ denote
    the corresponding triangulated subcategory of $\eK(\cA)$. When $\cA$
    is abelian, one may take for $O$ either the set $I$ of injectives or
    the set $P$ of projectives.
\end{itemize}

\subsubsection{}\label{VIII:1-2-4}

The category $\eK(\cA)$ is uniquely determined by the data of $\eC(\cA)$
and the family $\sC$. More precisely, every graded additive functor from
$\eC(\cA)$ to a triangulated category that takes every triangle of $\sC$
to a distinguished triangle factors uniquely through $\eK(\cA)$, and the
resulting functor is exact.

For every complex $X^\bullet$ of $\eC(\cA)$, let $\bar X^\bullet$ denote
the complex obtained from $X^\bullet$ by setting the differential equal
to zero. A three-term sequence
\[\xymatrix{
  0 \ar[r]
    & X^\bullet \ar[r]
    & Y^\bullet \ar[r]
    & Z^\bullet \ar[r]
    & 0
}\]
is called a semi-split sequence if the sequence
\[\xymatrix{
  0 \ar[r]
    & \bar X^\bullet \ar[r]
    & \bar Y^\bullet \ar[r]
    & \bar Z^\bullet \ar[r]
    & 0
}\]
is split; that is, if it is isomorphic to the sequence
\[\xymatrix{
  0 \ar[r]
    & \bar X^\bullet \ar[r]
    & \bar X^\bullet + \bar Z^\bullet \ar[r]
    & \bar Z^\bullet \ar[r]
    & 0 \text{.}
}\]

Let $0 \to X^\bullet \to Y^\bullet \to Z^\bullet \to 0$ be a semi-split
sequence. Choose a splitting, that is, for every $n$ an isomorphism
\[
  (Y^\bullet)^n \simeq (X^\bullet)^n + (Z^\bullet)^n \text{.}
\]
The matrix of the differential of $Y^\bullet$ in the above direct-sum
decomposition has the form
\[
  \begin{pmatrix}
    d_{X^\bullet} & \varphi_n \\
    0 & d_{Z^\bullet}
  \end{pmatrix}
\]
where $\varphi_n$ is a morphism from $(Z^\bullet)^n$ to
$(X^\bullet)^{n+1}$. The $\varphi_n$ determine a morphism of complexes
\[
  \varphi:Z^\bullet \to T(X^\bullet) \text{.}
\]
When the splitting is changed, the homotopy class of $\varphi$ does not
change. Moreover, if $\underline u$, $\underline v$, and
$\underline\varphi$ denote the images in $\eK(\cA)$ of the morphisms $u$,
$v$, and $\varphi$, then the triangle
$(X^\bullet,Y^\bullet,Z^\bullet,\underline u,\underline v,\underline\varphi)$
is distinguished. Let $\mathsf{SSS}(\cA)$ denote the category of
semi-split sequences of $\eC(\cA)$ and let $\mathsf{TrK}(\cA)$ denote the
category of distinguished triangles of $\eK(\cA)$. The preceding
construction defines a functor
\[
  \rho:\mathsf{SSS}(\cA) \to \mathsf{TrK}(\cA) \text{.}
\]
This functor is essentially surjective.

\subsection{Examples of cohomological and exact functors}\label{VIII:1-3}

\begin{definition}
Let $\cA$ be a triangulated category and let $\cB$ be an abelian category.
An additive functor $F:\cA\to \cB$ is called a \emph{cohomological
functor} if, for every distinguished triangle $(X,Y,Z,u,v,w)$, the
sequence
\[\xymatrix{
  F(X) \ar[r]^-{F(u)}
    & F(Y) \ar[r]^-{F(v)}
    & F(Z)
}\]
is exact.
\end{definition}

The functor $F\circ T^i$ will often be denoted by $F^i$. By axiom
\hyperlink{VIII:TR2}{TR2} for triangulated categories, there is an
unbounded exact sequence
\[\xymatrix{
  \cdots \ar[r]
    & F^i(X) \ar[r]
    & F^i(Y) \ar[r]
    & F^i(Z) \ar[r]
    & F^{i+1}(X) \ar[r]
    & \cdots
}\]

\subsubsection{Examples}\label{VIII:1-3-2}

\begin{enumerate}
  \item Let $\cA$ be a triangulated category. For every object $X$ of
    $\cA$, the functor $\hom_\cA(X,-)$ is cohomological. The functor
    $\hom_\cA(-,X)$ is a cohomological functor
    $\cA^\circ\to\mathsf{Ab}$.
  \item Let $\cA$ be an abelian category and let $\eC(\cA)$ be the category
    of complexes of $\cA$. Define the functor
    $\h^0:\eC(\cA)\to\cA$ as follows. For an object $X^\bullet$ of
    $\eC(\cA)$, set
    \[
      \h^0(X^\bullet) = \ker d^0 / \im{d^{-1}} \text{.}
    \]
    The functor $\h^0$ annihilates null-homotopic morphisms and therefore
    factors uniquely through $\eK(\cA)$. The resulting functor will again
    be denoted by $\h^0:\eK(\cA) \to \cA$.
\end{enumerate}

\subsubsection{Example of an exact bifunctor}\label{VIII:1-3-3}

Let $\cA$, $\cA'$, and $\cA''$ be three additive categories, and let
\[
  P:\cA\times \cA'\to \cA''
\]
be a bilinear functor (that is, additive in each argument separately).
It induces the bilinear functor
\[
  P^\ast:\eC(\cA)\times \eC(\cA') \to \eC(\cA'')
\]
as follows.

Let $X^\bullet$ be an object of $\eC(\cA)$ and let $Y^\bullet$ be an
object of $\eC(\cA')$. Then $P(X^\bullet,Y^\bullet)$ is a double complex
of $\cA''$. Set $P^\ast(X^\bullet,Y^\bullet) = $ the simple complex
associated with $P(X^\bullet,Y^\bullet)$.

Let $f$ be a null-homotopic morphism of $\eC(\cA)$ (respectively
$\eC(\cA')$), and let $Z^\bullet$ be an object of $\eC(\cA')$
(respectively $\eC(\cA)$). The morphism $P^\ast(f,Z^\bullet)$
(respectively $P^\ast(Z^\bullet,f)$) is then null-homotopic. It follows
that $P^\ast$ uniquely defines a functor
\[
  P^\bullet:\eK(\cA) \times \eK(\cA') \to \eK(\cA'') \text{.}
\]
The bifunctor $P^\bullet$ is exact.

In particular, let $\cA$ be an additive category. For $P$ one may take
the functor
\begin{align*}
  \cA^\circ\times \cA &\to \mathsf{Ab} \\
  (X,Y) &\mapsto \hom(X,Y)
\end{align*}

The preceding construction then gives a functor
\[
  \hom^\bullet:\eK(\cA)^\circ \times \eK(\cA) \to \eK(\mathsf{Ab})
\]
which, when composed with the functor
$\h^0:\eK(\mathsf{Ab})\to \mathsf{Ab}$, plainly gives back the functor
$\hom_{\eK(\cA)}$.

\section{Quotient categories}\label{VIII:2}

\subsection{Thick subcategories. Multiplicative systems of morphisms}\label{VIII:2-1}

We begin with two definitions.

\begin{definition}\label{VIII:2-1-1}
A subcategory $\cB$ of a triangulated category $\cA$ is said to be
\emph{thick} if $\cB$ is a full triangulated subcategory of $\cA$ and,
in addition, $\cB$ has the following property:

For every morphism $f:X\to Y$ that factors through an object of $\cB$ and
is contained in a distinguished triangle $(X,Y,Z,f,g,h)$ in which $Z$ is
an object of $\cB$, the source of $f$ and the target of $f$ are objects of
$\cB$.
\end{definition}

The set of thick subcategories of $\cA$ will be denoted by $\sN$.

The intersection of an arbitrary family of thick subcategories is a thick
subcategory.

\begin{definition}
Let $\cA$ be a triangulated category. A set $S$ of morphisms of $\cA$ is
called a \emph{multiplicative system of morphisms} if it has the following
five properties.

\paragraph{FR1}
\phantomsection\hypertarget{VIII:FR1}{}
If $f,g\in S$ and $f$ and $g$ are composable, then $f\circ g\in S$. For
every object $X$ of $\cA$, the identity morphism of $X$ belongs to $S$.

\paragraph{FR2}
\phantomsection\hypertarget{VIII:FR2}{}
In the category $\cA$, every diagram
\[\xymatrix{
  & Y \ar[d]^-{s\in S} \\
  Z \ar[r]^-f
    & X
}\]
can be completed to a commutative diagram
\[\xymatrix{
  P \ar[r]^-g \ar[d]^-{t\in S}
    & Y \ar[d]^-{s\in S} \\
  Z \ar[r]^-f
    & X
}\]
Moreover, the symmetric property holds.

\paragraph{FR3}
\phantomsection\hypertarget{VIII:FR3}{}
If $f$ and $g$ are morphisms, the following properties are equivalent:
\begin{enumerate}[\indent i)]
  \item there exists $s\in S$ such that $s\circ f=s\circ g$;
  \item there exists $t\in S$ such that $f\circ t=g\circ t$.
\end{enumerate}

\paragraph{FR4}
\phantomsection\hypertarget{VIII:FR4}{}
Let $T$ be the translation functor of $\cA$. For every element $s$ of $S$,
one has $T(s)\in S$.

\paragraph{FR5}
\phantomsection\hypertarget{VIII:FR5}{}
Let $(X,Y,Z,u,v,w)$ and $(X',Y',Z',u',v',w')$ be two distinguished
triangles, and let $f$ and $g$ be two elements of $S$ such that the pair
$(f,g)$ is a morphism from $u$ to $u'$. There exists a morphism $h\in S$
such that $(f,g,h)$ is a morphism of triangles.
\end{definition}

A multiplicative system of morphisms is said to be \emph{saturated} if it
has the following property:

A morphism $f$ belongs to $S$ if and only if there exist two morphisms $g$
and $g'$ such that $g\circ f\in S$ and $f\circ g'\in S$.

The intersection of an arbitrary family of saturated multiplicative
systems is a saturated multiplicative system.

The set of saturated multiplicative systems will be denoted by $\sS$.

Let $\cB$ be a thick subcategory. Denote by $\varphi(\cB)$ the set of
morphisms $f$ contained in a distinguished triangle $(X,Y,Z,f,g,h)$ in
which $Z$ is an object of $\cB$. Then $\varphi(\cB)$ is a saturated
multiplicative system.

Let $S$ be a saturated multiplicative system. Denote by $\psi(S)$ the
full subcategory generated by the objects $Z$ contained in a distinguished
triangle $(X,Y,Z,f,g,h)$ in which $f$ is an element of $S$. The
subcategory $\psi(S)$ is thick.

The main result of this subsection is the following: $\varphi$ is an
order-preserving isomorphism from $\sN$ onto $\sS$, where the order
relation is defined by inclusion. Its inverse is precisely $\psi$.

\subsection{Universal problems}\label{VIII:2-2}

Let $\cA$ and $\cA'$ be two triangulated categories and let $F$ be an exact
functor from $\cA$ to $\cA'$. Let $S(F)$ be the set of morphisms of $\cA$
that $F$ takes to isomorphisms, and let $\cB(F)$ be the full subcategory
generated by the objects of $\cA$ that $F$ takes to zero objects of
$\cA'$. Then $S(F)$ is a saturated multiplicative system and $\cB(F)$ is
a thick subcategory. Moreover, $S(F) = \varphi(\cB(F))$. Now let $\cA$ be
a triangulated category, let $\cB$ be a thick subcategory, and let
$S=\varphi(\cB)$ be the corresponding saturated multiplicative system.
The following two universal problems are equivalent.

\paragraph{Problem 1}
\phantomsection\hypertarget{VIII:prob1}{}
Find a triangulated category $\cA/\cB$ and an exact functor
$Q:\cA\to \cA/\cB$ such that every exact functor from $\cA$ to a
triangulated category $\cA'$ that takes the objects of $\cB$ to zero
objects of $\cA'$ admits a unique factorization of the form $G\circ Q$,
where $G$ is exact.

\paragraph{Problem 2}
\phantomsection\hypertarget{VIII:prob2}{}
Find a triangulated category $\cA_S$ and an exact functor
$Q:\cA\to \cA_S$ such that every exact functor from $\cA$ to a
triangulated category $\cA'$ that takes the morphisms of $S$ to
isomorphisms of $\cA'$ admits a unique factorization of the form
$G\circ Q$, where $G$ is exact.

(The author asks the reader to excuse this use of obsolete and reactionary
language.) In the next subsection we shall show that
\hyperlink{VIII:prob2}{Problem 2} has a solution. Consequently, the
corresponding \hyperlink{VIII:prob1}{Problem 1} has the same solution.

\subsection{Calculus of fractions}\label{VIII:2-3}

\subsubsection{}\label{VIII:2-3-1}

Let $\cA$ be a triangulated category and let $S$ be a saturated
multiplicative system. Denote by $\cA[S^{-1}]$ the category of fractions
of $\cA$ with respect to the set of morphisms $S$, and by $Q$ the
canonical functor from $\cA$ to $\cA[S^{-1}]$. The pair
$(\cA[S^{-1}],Q)$ solves the following problem: every functor from $\cA$
to an arbitrary category (not necessarily additive) that takes morphisms
of $S$ to isomorphisms factors uniquely through $(\cA[S^{-1}],Q)$. Such
a pair always exists, without any hypothesis on the set of morphisms $S$
\cite[1.1]{gz67}\footnote{The original text references [C.G.G] here. Most
likely, this is the lost work \emph{Catégories et foncteurs} by Chevalley,
Gabriel and Grothendieck, referenced in\cite{sc72}. }. However, because the
set of morphisms $S$ is a saturated multiplicative system, it has
properties \hyperlink{VIII:FR1}{FR1}, \hyperlink{VIII:FR2}{FR2}, and
\hyperlink{VIII:FR3}{FR3}. Consequently, by \cite{gz67}, the category
$\cA[S^{-1}]$ can be obtained by a right or left calculus of fractions.
We recall how $\cA[S^{-1}]$ is obtained by a right calculus of fractions.
(The left calculus of fractions is obtained by reversing the arrows.)

\subsubsection{}\label{VIII:2-3-2}

The objects of $\cA[S^{-1}]$ are the objects of $\cA$.
\begin{itemize}
  \item For every object $X$ of $\cA$, let $S_X$ denote the category of
    arrows belonging to $S$ and having target $X$. To every object $Y$ of
    $\cA$ one associates, functorially in $Y$, the following set-valued
    functor on $S_X^\circ$ ($S_X^\circ$ denotes the opposite category):
    \[
      H_Y:s\mapsto \hom(\operatorname{source}(s),Y) \text{.}
    \]
    Set
    \[
      \hom_{\cA[S^{-1}]}(X,Y) = \varinjlim_{S_X^\circ} H_Y \text{.}
    \]
    For all matters concerning direct limits, we refer to \cite{ar62}.

    One then observes with pleasure, using \hyperlink{VIII:FR1}{FR1},
    \hyperlink{VIII:FR2}{FR2}, and \hyperlink{VIII:FR3}{FR3}, that the
    category $S_X^\circ$ has properties L1, L2, and L3. Direct limits
    therefore have all the good properties of direct limits over filtered
    ordered sets.
  \item Let $X$, $Y$, and $Z$ be three objects of $\cA$, let $a$ be an
    element of $\hom_{\cA[S^{-1}]}(X,Y)$, and let $b$ be an element of
    $\hom_{\cA[S^{-1}]}(Y,Z)$. Let $s$ be an object of $S_X$, let
    $\underline a$ be an element of
    $\hom(\operatorname{source}(s),Y)$ whose image is $a$, and let an object
    of $S_Y$, $\underline b$ be an element of
    $\hom(\operatorname{source}(t),Z)$ whose image is $b$. There is then a
    diagram
    \begin{equation*}\tag{1}\label{VIII:eq:2-3-2-1}
    \xymatrix{
      & \bullet \ar[dl]_-s \ar[dr]^-{\underline a}
        &
        & \bullet \ar[dl]_-t \ar[dr]^-{\underline b} \\
      X
        &
        & Y
        &
        & Z
    }
    \end{equation*}
    which, by \hyperlink{VIII:FR2}{FR2}, can be completed to a diagram
    \[\xymatrix{
      & & \bullet \ar[dl]_-{t'} \ar[dr]^-{\underline c} \\
      & \bullet \ar[dl]_-s \ar[dr]^-{\underline a}
        & & \bullet \ar[dl]_-t \ar[dr]^-{\underline b} \\
      X
        & & Y
        & & Z
    }\]
    Let $b\circ a$ denote the image in $\hom_{\cA[S^{-1}]}(X,Z)$ of the
    morphism $\underline{b\circ c}$. Using properties
    \hyperlink{VIII:FR1}{FR1}, \hyperlink{VIII:FR2}{FR2}, and
    \hyperlink{VIII:FR3}{FR3}, one checks that $b\circ a$ is independent
    of the chosen representatives $\underline a$ and $\underline b$, and
    also of the way in which diagram \eqref{VIII:eq:2-3-2-1} is completed.
    One further checks that this defines a category $\cA[S^{-1}]$.
\end{itemize}

We denote by $Q$ the evident functor $\cA\to \cA[S^{-1}]$.

Since $\cA$ is additive, $\cA[S^{-1}]$ is an additive category. Moreover,
using \hyperlink{VIII:FR4}{FR4}, one proves that there exists a unique
translation functor on $\cA[S^{-1}]$, denoted by $T$, satisfying
\[
  Q\circ T = T\circ Q \text{.}
\]
Finally, using \hyperlink{VIII:FR5}{FR5}, one proves that there exists a
unique triangulated structure on $\cA[S^{-1}]$ for which $Q$ is exact.
The distinguished triangles of $\cA[S^{-1}]$ are precisely the triangles
isomorphic to the images under $Q$ of the distinguished triangles of
$\cA$. Denoting by $\cA_S$ the category $\cA[S^{-1}]$ equipped with this
triangulated structure, one proves without difficulty that the pair
$(\cA_S,Q)$ is a solution to \hyperlink{VIII:prob2}{Problem 2}.

\begin{definition}\label{VIII:2-3-3}
Let $\cA$ be a triangulated category, let $\cB$ be a thick subcategory of
$\cA$, and let $(Q,\cA/\cB)$ be the solution to
\hyperlink{VIII:prob1}{Problem 1} (\S \ref{VIII:2-2}). The category
$\cA/\cB$ is called the \emph{quotient category} of $\cA$ by $\cB$, and
$Q$ is called the canonical quotient functor. More generally, let $\cN$
be a triangulated subcategory of $\cA$, and let $\underline\cN$ be the
smallest thick subcategory containing $\cN$. By definition, the quotient
category of $\cA$ by $\cN$, denoted $\cA/\cN$, is
$\cA/\underline\cN$.
\end{definition}

\subsection{Properties of quotient categories}\label{VIII:2-4}

\subsubsection{}\label{VIII:2-4-1}

Let $\cA$ be a triangulated category and let $H$ be a cohomological functor
with values in an abelian category $\cG$. Let $\cB_H$ be the full
subcategory generated by the objects all of whose translates are taken by
$H$ to zero objects of $\cG$, and let $S_H$ be the set of morphisms all of
whose translates are taken by $H$ to isomorphisms of $\cG$. Then $\cB_H$
is a thick subcategory of $\cA$, and $S_H$ is a saturated multiplicative
system. Moreover, $S_H=\varphi(\cB_H)$ (\S \ref{VIII:2-1}). The functor
$H$ factors uniquely through $(Q,\cA/\cB_H)$.

\begin{theorem}\label{VIII:2-4-2}
Let $\cA$ be a triangulated category, let $\cB$ be a triangulated
subcategory, let $\cN$ be a thick subcategory of $\cA$, and let
$S=\varphi(\cN)$ be the corresponding saturated multiplicative system.
\begin{enumerate}[(a)]
  \item The category $\cN\cap\cB$ is a thick subcategory of $\cB$. The
    corresponding multiplicative system is $S\cap \cB$.
  \item The following two properties are equivalent:
    \begin{enumerate}[(i)]
      \item For every object $X$ of $\cB$ and every morphism $s:R\to X$,
        where $s$ belongs to $S$, there exists a morphism $s':R'\to R$ such
        that $s\circ s'$ belongs to $S\cap \cB$; and
        $R'\in\operatorname{Ob}\cB$.
      \item Every morphism from an object $X$ of $\cB$ to an object $Y$ of
        $\cN$ factors through an object of $\cN\cap \cB$.
    \end{enumerate}
  \item The properties (i)' and (ii)' obtained from (i) and (ii) by
    reversing the arrows are equivalent.
  \item If properties (i) and (ii) hold, or if properties (i)' and (ii)'
    hold, the canonical functor
    \[\xymatrix{
      \cB/\cN\cap \cB \ar[r] & \cA/\cN
    }\]
    is faithful.

    Since this functor is injective on objects, it realizes
    $\cB/\cN\cap\cB$ as a subcategory of $\cA/\cN$.
  \item If, in addition, $\cB$ is a full subcategory of $\cA$, then
    $\cB/\cN\cap \cB$ is a full subcategory of $\cA/\cN$.
\end{enumerate}
\end{theorem}

\begin{corollary}\label{VIII:2-4-3}
Let $\cA\subset \cB\subset \cC$ be three triangulated categories, with
$\cA$ a thick subcategory of $\cB$ and $\cB$ a thick subcategory of
$\cC$. Then $\cA$ is a thick subcategory of $\cC$, and $\cB/\cA$ is a
thick subcategory of $\cC/\cA$. The canonical functor
$\cC/\cB \to (\cC/\cA)/(\cB/\cA)$ is an isomorphism.
\end{corollary}

\subsection{Properties of the quotient functor}\label{VIII:2-5}

Let $\cA$ be a triangulated category, let $\cB$ be a thick subcategory,
let $S$ be the corresponding multiplicative system, and let
$Q:\cA\to \cA/\cB$ be the canonical quotient functor.

\begin{proposition}\label{VIII:2-5-1}
For objects $X$ and $Y$ of $\cA$, the following assertions are equivalent:
\begin{enumerate}[(a)]
  \item Every morphism $f:X\to Y$ for which there exists a morphism
    $s:Z\to X$ in $S$ satisfying $f\circ s=0$ is zero.
  \item Every morphism $f:X\to Y$ for which there exists a morphism
    $s:Y\to Z$, $s\in S$, satisfying $s\circ f=0$ is zero.
  \item Every morphism $f:X\to Y$ that factors through an object of $\cB$
    is zero.
  \item The canonical map
    \[
      \hom_\cA(X,Y) \to \hom_{\cA/\cB}(Q(X),Q(Y))
    \]
    is injective.
\end{enumerate}
\end{proposition}

\begin{proposition}\label{VIII:2-5-2}
Let $X$ and $Y$ be two objects of $\cA$. The following assertions are
equivalent:
\begin{enumerate}[(a)]
  \item Every diagram ($s\in S$)
    \[\xymatrix{
      & Z \ar[dl]_-s \ar[dr]^-f \\
      X
        & & Y
    }\]
    can be completed to a diagram: $(s,t\in S$):
    \[\xymatrix{
      & Z' \ar[d]_-t \\
      & Z \ar[dl]_-s \ar[dr]^-f \\
      X
        & & Y
    }\]
    in which $f\circ t$ factors through $(X,s\circ t)$.
  \item The assertion obtained by reversing the direction of the arrows in
    (a) and interchanging $X$ and $Y$.
  \item Every diagram
    \[\xymatrix{
      X \ar[r]^-f
        & N \ar[r]^-g
        & Y
    }\]
    in which $g\circ f=0$ and $N$ is an object of $\cB$ can be completed
    to a diagram
    \[\xymatrix{
      & N' \ar[d]^-i \\
      X \ar[ur]^-h \ar[r]^-f
        & N \ar[r]^-g
        & Y
    }\]
    in which $f=i\circ h$ and $g\circ i=0$.
  \item The assertion obtained by reversing the arrows in (c) and
    interchanging $X$ and $Y$.
  \item The canonical map
    \[\xymatrix{
      \hom_\cA(X,Y) \ar[r]^-Q
        & \hom_{\cA/\cB}(Q(X),Q(Y))
    }\]
    is surjective.
\end{enumerate}
\end{proposition}

\begin{proposition}\label{VIII:2-5-3}
Let $X$ be an object of $\cA$. The following assertions are equivalent:
\begin{enumerate}[(a)]
  \item For every object $Y$ of $\cA$, the canonical map
    \[\xymatrix{
      \hom_\cA(X,Y) \ar[r]^-Q
        & \hom_{\cA/\cB}(Q(X),Q(Y))
    }\]
    is an isomorphism.
  \item Every morphism from $X$ to an object of $\cB$ is zero.
\end{enumerate}
Similarly, the following two assertions are equivalent:
\begin{enumerate}[(a)']
  \item For every object $Y$ of $\cA$, the canonical map
    \[\xymatrix{
      \hom_\cA(Y,X) \ar[r]^-Q
        & \hom_{\cA/\cB}(Q(Y),Q(X))
    }\]
    is an isomorphism.
  \item Every morphism from an object of $\cB$ to $X$ is zero.
\end{enumerate}
\end{proposition}

\begin{definition}\label{VIII:2-5-4}
Every object having the equivalent properties (a) and (b) of the preceding
proposition is called an \emph{object free on the left over $\cA/\cB$},
or a \emph{left $Q$-free object}. Similarly, every object having the
equivalent properties (a)' and (b)' is called an \emph{object free on the
right over $\cA/\cB$}, or a \emph{right $Q$-free object}. It is determined
up to isomorphism by the knowledge of
$Q(X)\in\operatorname{Ob}{\cA/\cB}$.
\end{definition}

\subsection{Functors adjoint to the quotient functor}\label{VIII:2-6}

\begin{definition}\label{VIII:2-6-1}
Two triangulated subcategories $\cN$ and $\cN'$ of a triangulated category
$\cA$ are said to be \emph{orthogonal} if, for every object $X$ of $\cN$
and every object $Y$ of $\cN'$, one has
\[
  \hom_\cA(X,Y) = 0
\]
The category $\cN'$ is then said to be right orthogonal to $\cN$, and
$\cN$ left orthogonal to $\cN'$.
\end{definition}

\begin{proposition}\label{VIII:2-6-2}
Let $\cN$ be a triangulated subcategory of a triangulated category $\cA$.
The full category $\cN^\bot$ (respectively $^\bot\cN$) generated by the
objects $X$ of $\cA$ such that, for every object $Y$ of $\cN$,
$\hom_\cA(Y,X)=0$ (respectively $\hom_\cA(X,Y)=0$), is a thick subcategory
of $\cA$. By an abuse of language, the category $\cN^\bot$ is called the
right orthogonal of $\cN$.
\end{proposition}

This proposition allows us to restate Proposition \ref{VIII:2-5-3}.

\begin{proposition}\label{VIII:2-6-3}
Let $\cB\subset\cA$ be a thick subcategory of a triangulated category
$\cA$. The full category generated by the objects free on the right over
$\cA/\cB$ (respectively on the left) is precisely the right
(respectively left) orthogonal of $\cB$.
\end{proposition}

\begin{proposition}\label{VIII:2-6-4}
Let $\cA$ be a triangulated category, let $\cB$ be a thick subcategory of
$\cA$, let $\cB^\bot$ be the right orthogonal of $\cB$, let $S(\cB)$ and
$S(\cB^\bot)$ be the corresponding multiplicative systems, and let $Q$
and $Q^\bot$ be the canonical quotient functors. For an object $X$ of
$\cA$, consider the following properties:
\begin{enumerate}[(i)]
  \item $X$ maps by a morphism of $S(\cB)$ to an object of $\cB^\bot$.
  \item $X$ receives, by a morphism of $S(\cB^\bot)$, an object of
    $\cB$.
  \item The category $S_X(\cB)$ of arrows of $S(\cB)$ with source $X$ has
    a final object.
  \item The category $S^X(\cB^\bot)$ of arrows of $S(\cB^\bot)$ with
    target $X$ has an initial object.
  \item The category $\cB/X$ of objects of $\cB$ over $X$ has a final
    object.
  \item The category $\cB^\bot/X$ of objects of $\cB^\bot$ under $X$ has
    an initial object.
\end{enumerate}
Then
\[\xymatrix@=0.5cm{
  & & & & (iv) \\
  (i) \ar@{<=>}[r]
    & (ii) \ar@{<=>}[r]
    & (iii) \ar@{<=>}[r]
    & (v) \ar@{=>}[ur] \ar@{=>}[dr] \\
  & & & & (vi)
}\]

If, in addition, $^\bot(\cB^\bot)=\cB$, then all the properties are
equivalent.
\end{proposition}

Let $\cB$ be a thick subcategory of a triangulated category $\cA$, let
$i$ be the inclusion functor from $\cB$ into $\cA$, let $Q$ be the
quotient functor from $\cA$ to $\cA/\cB$, let $\cB^\bot$ be the right
orthogonal of $\cB$, and let $i^\bot$ and $Q^\bot$ be the corresponding
functors. Proposition \ref{VIII:2-6-4} immediately gives the following
propositions.

\begin{proposition}\label{VIII:2-6-5}
The following two properties are equivalent:
\begin{enumerate}[(i)]
  \item The functor $i$ has a right adjoint.
  \item The functor $Q$ has a right adjoint.
\end{enumerate}
\end{proposition}

The proposition remains true when “right” is replaced by “left.”

\begin{proposition}\label{VIII:2-6-6}
The following two properties are equivalent:
\begin{enumerate}[(i)]
  \item The functor $i$ has a right adjoint.
  \item The functor $i^\bot$ has a left adjoint. The left orthogonal of the
    category $\cB^\bot$ is equal to $\cB$.
\end{enumerate}
\end{proposition}

\begin{proposition}\label{VIII:2-6-7}
Suppose the properties of Propositions \ref{VIII:2-6-5} and
\ref{VIII:2-6-6} hold. Let $i^\ast$ and $Q^\ast$ be the right adjoints of
$i$ and $Q$. Similarly, let $^\ast i^\bot$ and $^\ast Q^\bot$ be the left
adjoints of $i^\bot$ and $Q^\bot$. All these functors are exact. Moreover:

There is a functorial isomorphism
$Q^\ast\circ Q\iso i^\bot\circ ^\ast i^\bot$ such that the following
diagram is commutative:
\[\xymatrix{
  & \operatorname{id} \ar[dl] \ar[dr] \\
  Q^\ast \circ Q \ar[rr]^\sim
    & & i^\bot\circ ^\ast i^\bot
}\]
(the oblique arrows are defined by the adjunction properties).

Similarly, there is an isomorphism
$^\ast Q^\bot \circ Q^\bot \iso i\circ i^\ast$ such that the following
diagram is commutative:
\[\xymatrix{
  ^\ast Q^\bot\circ Q^\bot \ar[rr]^-\sim \ar[dr]
    & & i\circ i^\ast \ar[dl] \\
  & \operatorname{id}{}
}\]

Finally, there is a functorial morphism of degree $1$,
$\partial:i^\bot\circ ^\ast i^\bot \to i\circ i^\ast$, such that the
following triangle is distinguished $(\deg\partial=1$):
\[\xymatrix{
  i^\bot \circ ^\ast i^\bot \ar[rr]^-\partial
    & & i\circ i^\ast \ar[dl] \\
  & \operatorname{id}{} \ar[ul]
}\]

Such a morphism $\partial$ is unique and satisfies
$\partial(TX)=-T\partial(X)$ ($X\in\operatorname{Ob}\cA$, $T$ the
translation).
\end{proposition}

\section{Derived categories of an abelian category}\label{VIII:3}

\subsection{Derived categories}\label{VIII:3-1}

We use the notation of \ref{VIII:1-2-3}. We shall also use the following
notation.

\subsubsection{Notation}\label{VIII:3-1-1}

Let $\cA$ be an abelian category. Set
\begin{align*}
  \eD(\cA) &= \eK^{\infty,\infty}(\cA)/\eK^{\infty,\varnothing}(\cA) \\
  \eD^b(\cA) &= \eK^{b,b}(\cA)/\eK^{b,\varnothing}(\cA) \\
  \eD^+(\cA) &= \eK^{+,+}(\cA) / \eK^{+,\varnothing}(\cA) \\
  \eD^-(\cA) &= \eK^{-,-}(\cA) / \eK^{-,\varnothing}(\cA) \text{.}
\end{align*}

\begin{proposition}\label{VIII:3-1-2}
The natural functors appearing in the following diagram
\[\xymatrix@=0.5cm{
  & \eD(\cA) \\
  \eD^+(\cA) \ar[ur]
    & & \eD^-(\cA) \ar[ul] \\
  & \eD^b(\cA) \ar[ul] \ar[ur]
}\]
are fully faithful.

The natural functors
\begin{align*}
  \eD^b(\cA) &= \eK^{b,b}(\cA)/\eK^{b,\varnothing}(\cA) \to \eK^{+,b}(\cA)/\eK^{+,\varnothing}(\cA) \\
  \eD^b(\cA) &= \eK^{b,b}(\cA)/\eK^{b,\varnothing}(\cA) \to \eK^{-,b}(\cA) / \eK^{-,\varnothing}(\cA)
\end{align*}
are equivalences of categories.

Since the functors in the diagram
\[\xymatrix@=0.5cm{
  & \eD(\cA) \\
  \eD^+(\cA) \ar[ur]
    & & \eD^-(\cA) \ar[ul] \\
  & \eD^b(\cA) \ar[ul] \ar[ur]
}\]
are injective on objects, they realize the categories $\eD^b(\cA)$,
$\eD^+(\cA)$, and $\eD^-(\cA)$ as full subcategories of $\eD(\cA)$.
\end{proposition}

The proof of this proposition is contained in Theorem \ref{VIII:2-4-2}.

\begin{definition}\label{VIII:3-1-3}
The category $\eD(\cA)$ is called the \emph{derived category} of $\cA$.
The objects of $\eD(\cA)$ isomorphic to objects of $\eD^b(\cA)$ are
called the \emph{bounded} objects of $\eD(\cA)$. The objects of
$\eD(\cA)$ isomorphic to objects of $\eD^+(\cA)$ are called the
\emph{left-bounded} objects, and the objects of $\eD(\cA)$ isomorphic to
objects of $\eD^-(\cA)$ are called the \emph{right-bounded} objects.
\end{definition}

Let $D$ denote the canonical functor
\[
  D : \eC(\cA) \to \eD(\cA) \text{.}
\]
The “restriction” of $D$ to the full subcategory $\cA$ of $\eC(\cA)$
will again be denoted by $D$. The functor $D$ restricted to $\cA$ is fully
faithful. The functor $D$ factors uniquely through $\eK(\cA)$. Let $X$
and $Y$ be two objects of $\cA$.

One readily checks that, for every $m>0$, the groups
\[
  \hom_{\eD(\cA)}(D(X),T^{-m} D(Y)) \qquad \text{($T$ is the translation functor)}
\]
are zero. The cohomological dimension of $\cA$ is the least integer $n$
such that, for every $m>n$,
\[
  \hom_{\eD(\cA)}(D(X),T^m D(Y)) = 0 \qquad\text{for every pair $X,Y$ of objects of $\cA$.}
\]
If there is no such integer, the cohomological dimension of $\cA$ is said
to be infinite.

Let $\eI^+(\cA)$ (respectively $\eI(\cA)$) denote the full triangulated
subcategory of $\eK^+(\cA)$ (respectively of $\eK(\cA)$) defined by the
complexes whose objects are injective in every degree.

\begin{proposition}\label{VIII:3-1-4}
Suppose that the category $\cA$ has enough injectives. The natural functor
\[
  Q^+:\eK^+(\cA) \to \eD^+(\cA)
\]
induces an equivalence of categories
\[
  \eI^+(\cA) \to \eD^+(\cA)  \text{.}
\]
The quasi-inverse functor is a right adjoint of $Q^+$.
\end{proposition}

If, in addition, $\cA$ has finite cohomological dimension, the preceding
assertion remains true after the superscripts $+$ are omitted.

There is an analogous statement concerning projectives.

The proof of Proposition \ref{VIII:3-1-4} relies on
\ref{VIII:2-6-4}.

The triangulated structure of $\eD(\cA)$ is made more explicit by the
following proposition.

Let $\eE(\cA)$ be the category of three-term exact sequences in
$\eC(\cA)$. The category $\mathsf{SSS}(\cA)$ of semi-split exact
sequences (\S\ref{VIII:1-2-4}) is a full subcategory of $\eE(\cA)$.
We defined (loc. cit.) a functor
\[
  \rho:\mathsf{SSS}(\cA) \to \mathsf{TrK}(\cA)
\]
where $\mathsf{TrK}(\cA)$ denotes the category of distinguished triangles
of $\eK(\cA)$. This gives a functor
\[
  \sigma : \mathsf{SSS}(\cA) \to \mathsf{TrD}(\cA)
\]
where $\mathsf{TrD}(\cA)$ denotes the category of distinguished triangles
of $\eD(\cA)$.

\begin{proposition}\label{VIII:3-1-5}
There exists a unique functor
\[
  \Pi:\eE(\cA) \to \mathsf{TrD}(\cA)
\]
of the following form ($\deg \delta(S)=1$):
\[
(S=0 \to X^\bullet \to Y^\bullet \to Z^\bullet \to 0)
\mapsto
\xymatrix{
  & D(Z^\bullet) \ar[dl]_-{\delta(S)} \\
  D(X^\bullet) \ar[rr]^-{D(u)}
    & & D(Y^\bullet) \ar[ul]_-{D(v)}
}
\]
whose restriction to $\mathsf{SSS}(\cA)$ is $\sigma$. \emph{This functor
is essentially surjective.}
\end{proposition}

\begin{proposition}\label{VIII:3-1-6}
The canonical cohomological functor
\[
  \h^0:\eK(\cA) \to \cA
\]
factors uniquely through a functor that we again denote by $\h^0$,
\[
  \h^0:\eD(\cA) \to \cA \text{.}
\]
On $\eD(\cA)$, the functor $\h^0$ and its translates $\h^i$ form a
conservative system; that is, a morphism $f:X^\bullet \to Y^\bullet$ is
an isomorphism if and only if, for every integer $i$,
\[
  \h^i(f):\h^i(X) \to \h^i(Y)
\]
is an isomorphism.
\end{proposition}

\begin{definition}\label{VIII:3-1-7}
A morphism $f:X^\bullet \to Y^\bullet$ of $\eC(\cA)$ (respectively of
$\eK(\cA)$) is called a \emph{quasi-isomorphism} if, for every integer
$i$,
\[
  \h^i(f):\h^i(X) \to \h^i(Y)
\]
is an isomorphism.
\end{definition}

By the preceding proposition, the quasi-isomorphisms are the morphisms
that become isomorphisms in $\eD(\cA)$.

\begin{definition}\label{VIII:3-1-8}
Let $\sigma$ be a set of objects of $\cA$ and let $X^\bullet$ be an
object of $\eC(\cA)$ (respectively $\eK(\cA)$). A right (respectively
left) resolution of type $\sigma$ of $X^\bullet$ is a quasi-isomorphism
$X^\bullet \to V^\bullet$ (respectively $V^\bullet \to X^\bullet$), where
$V^\bullet$ is a complex all of whose objects belong to $\sigma$. We say
\emph{injective resolution} (respectively projective resolution) instead
of right (respectively left) resolution of injective (respectively
projective) type.
\end{definition}

In general, we shall omit the words “right” or “left” when no ambiguity
can result.

\begin{proposition}\label{VIII:3-1-9}
\begin{enumerate}
  \item Suppose that every object of $\cA$ embeds in an object of $\sigma$
    (respectively is a quotient of an object of $\sigma$). Then every
    object of $\eC^+(\cA)$ (respectively $\eC^-(\cA)$) has a right
    (respectively left) resolution of type $\sigma$ in $\eC^+(\cA)$
    (respectively in $\eC^-(\cA)$).
  \item Suppose, in addition, that there exists an integer $n$ such that
    for every exact sequence
    \begin{align*}
      Y^0 &\to Y^1 \to \cdots \to Y^{n-1} \to Y^n \to 0 \\
      \text{(respectively } 0 &\to Y^n \to Y^{n-1} \to \cdots \to Y^1 \to Y^0 \text{)}
    \end{align*}
    of objects of $\cA$ in which, for every $i$, $0\leqslant i\leqslant n-1$,
    $Y^i$ is an object of $\sigma$, the object $Y^n$ is also an object of
    $\sigma$. Then every object of $\eC(\cA)$ has a right (respectively
    left) resolution of type $\sigma$.
\end{enumerate}
\end{proposition}

\subsection{Study of $\ext$}\label{VIII:3-2}

Let $X^\bullet$ and $Y^\bullet$ be two objects of $\eC(\cA)$
(respectively $\eK(\cA)$), and let
$D:\eC(\cA) \to \eD(\cA)$ (respectively
$Q:\eK(\cA) \to \eD(\cA)$) be the canonical functor.

\begin{definition}\label{VIII:3-2-1}
We call the \emph{$i$th hyper-Ext}, and denote by
\[
  (X^\bullet,Y^\bullet)\mapsto \ext^i(X^\bullet,Y^\bullet)
\]
the functor $\hom_{\eD(\cA)}(D(X^\bullet),T^i D(Y^\bullet))$
(respectively
$\hom_{\eD(\cA)}(Q(X^\bullet),T^i Q(Y^\bullet))$).
\end{definition}

\begin{proposition}\label{VIII:3-2-2}
Let $0 \to X^\bullet \to Y^\bullet \to Z^\bullet \to 0$ be an exact
sequence in $\eC(\cA)$, and let $V^\bullet$ be an object of $\eC(\cA)$.
There are unbounded exact sequences
\[\xymatrix@=0.5cm{
  \cdots \ar[r]
    & \ext^i(V^\bullet,X^\bullet) \ar[r]
    & \ext^i(V^\bullet,Y^\bullet) \ar[r]
    & \ext^i(V^\bullet,Z^\bullet) \ar[r]^-\delta
    & \ext^{i+1}(V^\bullet,X^\bullet) \ar[r]
    & \cdots \\
  \cdots \ar[r]
    & \ext^i(Z^\bullet,V^\bullet) \ar[r]
    & \ext^i(Y^\bullet,V^\bullet) \ar[r]
    & \ext^i(X^\bullet,V^\bullet) \ar[r]^-\delta
    & \ext^{i+1}(Z^\bullet,V^\bullet) \ar[r]
    & \cdots
}\]
\end{proposition}

Let $X^\bullet$ and $Y^\bullet$ be two objects of $\eK(\cA)$. Denote by
$\mathsf{Qis}/X^\bullet$ (respectively $\mathsf{Qis}^+/X^\bullet$,
$\mathsf{Qis}^-/X^\bullet$, $\mathsf{Qis}^b/X^\bullet$) the category of
quasi-isomorphisms with target $X^\bullet$ and source in $\eK(\cA)$
(respectively $\eK^+(\cA)$, $\eK^-(\cA)$, $\eK^b(\cA)$). Similarly,
define the categories $Y^\bullet/\mathsf{Qis}$,
$Y^\bullet/\mathsf{Qis}^+$,\ldots (quasi-isomorphisms with source
$Y^\bullet)$.

The following proposition summarizes the various equivalent definitions
of $\ext^i$.

\begin{proposition}\label{VIII:3-2-3}
1. There are isomorphisms of functors
\[\xymatrix{
  \ext^0(X^\bullet,Y^\bullet) = \hom_{\eD(\cA)}(Q(X^\bullet),Q(Y^\bullet)) \ar[r]^-\sim
    & \varinjlim_{(\mathsf{Qis}/X^\bullet)^\circ} \hom_{\eK(\cA)}(-,Y^\bullet) \ar[r]^-\sim
    & \\
  \varinjlim_{Y^\bullet/\mathsf{Qis}} \hom_{\eK(\cA)}(X^\bullet,-) \ar[r]^-\sim
    & \varinjlim_{(\mathsf{Qis}/X^\bullet)^\circ \times Y^\bullet/\mathsf{Qis}} \hom_{\eK(\cA)}(-,-) \text{.}
}\]
2. If $X^\bullet$ is an object of $\eK^-(\cA)$,
\[
  \ext^0(X^\bullet,Y^\bullet) \iso \varinjlim_{(\mathsf{Qis}^-/X^\bullet)^\circ} \hom_{\eK(\cA)}(-,Y^\bullet)
\]
If $Y^\bullet$ is an object of $\eK^+(\cA)$,
\[
  \ext^0(X^\bullet,Y^\bullet) \iso \varinjlim_{Y^\bullet/\mathsf{Qis}} \hom_{\eK(\cA)}(X^\bullet,-)
\]
3. If $X^\bullet$ is an object of $\eK^-(\cA)$ and $Y^\bullet$ is an object
of $\eK^b(\cA)$,
\[
  \ext^0(X^\bullet,Y^\bullet) \iso \varinjlim_{Y^\bullet/\mathsf{Qis}^b} \hom_{\eK(\cA)}(X^\bullet,-)
\]
If $X^\bullet$ is an object of $\eK^b(\cA)$ and $Y^\bullet$ is an object
of $\eK^+(\cA)$,
\[
  \ext^0(X^\bullet,Y^\bullet) \iso \varinjlim_{(\mathsf{Qis}^b/X^\bullet)^\circ} \hom_{\eK(\cA)}(-,Y^\bullet)
\]
4. If the category $\cA$ has enough injectives and $Y^\bullet$ is an object
of $\eK^+(\cA)$, then $Y^\bullet$ has an injective resolution, unique up
to isomorphism in $\eK^+(\cA)$,
\[
  Y^\bullet \to I(Y^\bullet)
\]
and
\[
  \ext^0(X^\bullet,Y^\bullet) \iso \hom_{\eK(\cA)}(X^\bullet,I(Y^\bullet)) \text{.}
\]
There is an analogous statement for projectives.
5. If the category $\cA$
has enough injectives and if $\cA$ has \emph{finite cohomological dimension}, then
every object $Y^\bullet$ of $\eK(\cA)$ has a unique injective resolution
\[
  Y^\bullet \to I(Y^\bullet)
\]
and
\[
  \ext^0(X^\bullet,Y^\bullet) \iso \hom_{\eK(\cA)}(X^\bullet,I(Y^\bullet))
\]
There is an analogous statement for projectives.
\end{proposition}

\paragraph{Remark:}
On making explicit the isomorphism of Proposition
\ref{VIII:3-2-3}(3), one recovers the Yoneda definition.

\section{Derived functors}\label{VIII:4}

\subsection{Definition of derived functors}\label{VIII:4-1}

\begin{definition}\label{VIII:4-1-1}
Let $\cC$ and $\cC'$ be two graded categories (we denote by $T$ the
translation functor of $\cC$ and $\cC'$), and let $F$ and $G$ be two
graded functors from $\cC$ to $\cC'$. A \emph{morphism of graded functors}
is a morphism of functors
\[
  u:F\to G
\]
having the following property:

For every object $X$ of $\cC$, the following diagram is commutative:
\[\xymatrix{
  F(T X) \ar[r]^-{u(T X)}
    & G(T X) \\
  T F(X) \ar[r]^-{T u(X)} \ar[u]_-\wr
    & T G(X) \ar[u]_-\wr
}\]
\end{definition}

Let $\cC$ and $\cC'$ be two triangulated categories. Denote by
$\mathsf{Fex}(\cC,\cC')$ the category of exact functors from $\cC$ to
$\cC'$, the morphisms between two functors being morphisms of graded
functors.

Let $\cA$ and $\cB$ be two abelian categories and let
$\Phi:\eK^\ast(\cA)\to \eK^{\ast'}(\cB)$ be an exact functor ($\ast$ and
$\ast'$ denote one of the signs $+$, $-$, $b$, or $v$, “empty”). By
composition, the canonical functor
$Q:\eK^\ast(\cA)\to \eD^\ast(\cA)$ gives a functor
\[\xymatrix{
  \mathsf{Fex}(\eD^\ast(\cA),\eD^{\ast'}(\cB)) \ar[r]
    & \mathsf{Fex}(\eK^\ast(\cA),\eD^{\ast'}(\cB))
}\]
and hence (also denoting by $Q'$ the canonical functor
$\eK^{\ast'}(\cB) \to \eD^{\ast'}(\cB)$) a functor $\%$
(respectively $\%'$)
$\mathsf{Fex}(\eD^\ast(\cA),\eD^{\ast'}(\cB)) \to \mathsf{Ab}$:
\begin{align*}
  \Psi  &\mapsto \hom(Q'\circ \Phi,\Psi\circ Q) \\
  (\text{respectively }\Psi &\mapsto \hom(\Psi\circ Q,Q'\circ \Phi)\text{ ).}
\end{align*}

\begin{definition}\label{VIII:4-1-2}
We say that $\Phi$ admits a \emph{total right derived functor}
(respectively total left derived functor) if the functor $\%$
(respectively $\%'$) is representable. An object representing $\%$
(respectively $\%'$) is called the total right (respectively left) derived
functor of $\Phi$ and is denoted by $\eR \Phi$ (respectively
$\eL \Phi$).\footnote{A more manageable definition (and one equivalent in
practice) is given in \cite[XVII 1.2]{sga4}.}
\end{definition}

\subsubsection{Notation}\label{VIII:4-1-3}

Let $\cA$ and $\cB$ be two abelian categories and let $f:\cA\to \cB$ be
an additive functor. The \emph{total right derived functor} of
$\eK^+(f):\eK^+(\cA) \to \eK^+(\cB)$ is denoted by $\eR^+ f$. The
functors $\eR^- f$, $\eR f$, $\eL^+ f$, $\eL^- f$, and $\eL f$ are
defined similarly.

Likewise, let $F:\eK^\ast(\cA) \to \cB$ be a cohomological functor.
By composition, the canonical quotient functor
$Q:\eK^\ast(\cA) \to \eD^\ast(\cA)$ defines a functor
\[\xymatrix{
  \mathsf{Foco}(\eD^\ast(\cA),\cB) \ar[r]
    & \mathsf{Foco}(\eK^\ast(\cA),\cB)
}\]
(where $\mathsf{Foco}(-,-)$ denotes the category of cohomological
functors), and hence a functor $\%$ (respectively $\%'$)
$\mathsf{Foco}(\eD^\ast(\cA),\cB) \to \mathsf{Ab}$ given by
$G\mapsto \hom(F,G\circ Q)$ (respectively
$G\mapsto \hom(G\circ Q,F)$).

\begin{definition}\label{VIII:4-1-4}
We say that the functor $F$ admits a \emph{right cohomological derived
functor} (respectively left cohomological derived functor) if the functor
$\%$ (respectively $\%'$) is representable. An object representing $\%$
(respectively $\%'$) is called the right (respectively left)
cohomological derived functor and is denoted by $\eR F$ (respectively
$\eL F$).
\end{definition}

\subsubsection{Notation}\label{VIII:4-1-5}

Let $\cA$ and $\cB$ be two abelian categories and let $f:\cA\to \cB$ be
an additive functor. The right cohomological derived functor of
\[
  \h^0{}\circ \eK^+ f:\eK^+(\cA) \to \cB
\]
is denoted by $\eR^+ f$. When no confusion can result, the functor
$\eR^+ f\circ T^i$ ($T$ is the translation functor) is denoted by
$\eR^i f$. The right cohomological derived functor of
$\h{}\circ \eK f$ is denoted by $\eR f$. When no confusion can result,
the functor $\eR f\circ T^i$ is denoted by $\eR^i f$.

The functors $\eL^- f$, $\eL f$, and $\eL^i f$ are defined similarly.

\subsubsection{Remarks}\label{VIII:4-1-6}

As we shall see in the next subsection, these definitions include the
classical definition of derived functors when the categories have enough
injectives or projectives and the complexes are suitably bounded. They are
nevertheless more general, but without further hypotheses they are not
very manageable. For example, let $f:\cA\to \cB$ be an additive functor
between two abelian categories. Suppose that the functors $\eR^+ f$ and
$\eR f$ exist. I do not know how to prove that $\eR f$ induces on
$\eD^+(\cA)$ the functor $\eR^+ f$. Suppose, in addition, that the functor
$\eR^+ f$ (cohomological derived) exists. I do not know how to prove that
the functor $\h^0{}\circ \eR^+ f$ is isomorphic to $\eR^+ f$. In practice,
however, these properties will hold.

\subsection{Existence of derived functors}\label{VIII:4-2}

\begin{proposition}\label{VIII:4-2-1}
The hypotheses are those of Definition \ref{VIII:4-1-4}. Suppose, in
addition, that $\cA$ is a $\sU$-category, where $\sU$ is a universe such
that the set $\dZ$ of rational integers belongs to $\sU$, and such that
the set of objects of $\cA$ belongs to $\sU$. Suppose also that $\cB$ is
the category of abelian $\sU$-groups, or more generally the category of
sheaves of abelian $\sU$-groups on an arbitrary $\sU$-site. (A $\sU$-site
is a $\sU$-category whose set of objects belongs to $\sU$, equipped with
a topology.) The functor $F$ admits a right cohomological derived functor.
This functor is calculated as follows:

Let $X$ be an object of $\eD^\ast(\cA)$, and let $\underline X$ be the
object of $\eK^\ast(\cA)$ above $X$. There is a functorial isomorphism
\[
  \eR F(X) = \varinjlim_{\underline X/\mathsf{Qis}^\ast} F(-) \text{.}
\]
\end{proposition}

In particular, let $f:\cA\to \cB$ be an additive functor. The functors
$\eR f$ and $\eR^+ f$ exist and are calculated as above. The functor
$\eR f$ induces on $\eD^+(\cA)$ the functor $\eR^+ f$. Since no confusion
can then result, we shall use the notation $\eR^i f$.

\subsubsection{Remarks}\label{VIII:4-2-2}

Proposition \ref{VIII:4-2-1} can be expressed in a more general form in
the setting of triangulated categories. The hypotheses on the category
$\cB$ needed to ensure the validity of Proposition \ref{VIII:4-2-1} are
hypotheses asserting the existence and exactness of direct limits over
pseudo-filtered $\sU$-categories whose sets of objects belong to $\sU$.
Proposition \ref{VIII:4-2-1} introduces an asymmetry between right and
left derived functors, at least in practice.

\begin{theorem}\label{VIII:4-2-3} % labelled 4.2.2 in original
The data are those of Definition \ref{VIII:4-1-2}. Suppose, in addition,
that a full triangulated subcategory of $\eK^\ast(\cA)$, denoted $\cC$, is given,
and denote by $\cN$ the full triangulated subcategory of the acyclic
objects of $\cC$. Suppose that
\begin{enumerate}[(1)]
  \item Every object of $\cN$ is taken by $\Phi$ to an acyclic object of
    $\eK^{\ast'}(\cB)$.
  \item Every object $X$ of $\eK^\ast(\cA)$ is the source (respectively
    target) of a quasi-isomorphism whose target (respectively source) is
    an object of $\cC$.
\end{enumerate}
\begin{enumerate}[a)]
  \item Then $\Phi$ admits a total right (respectively left) derived
    functor.
  \item The functor $\eR \Phi$ (respectively $\eL \Phi$) is obtained as
    follows:

    The restriction of the functor
    $Q'\circ \Phi:\eK^\ast(\cA) \to \eD^{\ast'}(\cB)$ to $\cC$ vanishes on the
    objects of $\cN$ and therefore factors through $\cC/\cN$. This gives
    a functor $\Phi':\cC/\cN\to \eD^{\ast'}(\cB)$. But the natural
    functor $\cC/\cN\to \eD^\ast(\cA)$ is an equivalence of categories
    (Theorem \ref{VIII:2-4-2}). Composing with its quasi-inverse gives the
    functor $\eR\Phi$ (respectively $\eL \Phi$).
  \item Let $X$ be an object of $\eK^\ast(\cA)$, let $Y$ be an object of
    $\cC$, and let $X\to Y$ (respectively $Y\to X$) be a
    quasi-isomorphism. The object $\eR\Phi\circ Q(X)$ (respectively
    $\eL \Phi\circ Q(X)$) is isomorphic to $Q'\circ \Phi(Y)$.
  \item The functor $X\mapsto \eR\Phi\circ Q(X)$ (respectively
    $X\mapsto \eL \Phi\circ Q(X)$) is isomorphic to the functor
    $X\mapsto \varinjlim_{X/\mathsf{Qis}^\ast} Q'\circ\Phi(-)$
    (respectively
    $X\mapsto \varinjlim_{\mathsf{Qis}^\ast/X} Q'\circ \Phi(-)$).
  \item The functor $\h^0{}\circ \Phi:\eK^\ast(\cA)\to \cB$ admits a
    right (respectively left) cohomological derived functor, which is
    precisely $\h^0{}\circ\eR\Phi$ (respectively
    $\h^0{}\circ \eL\Phi$).
\end{enumerate}
\end{theorem}

\begin{corollary}\label{VIII:4-2-3-1} % originally corollary 1
Suppose that $\ast=+$ (respectively $\ast=-$) in the hypotheses of the
preceding theorem and that $\cA$ has enough injectives (respectively
projectives). Then Theorem \ref{VIII:4-2-3} applies with $\cC$ the
category of complexes of injectives (respectively projectives). Suppose
in addition that $\cA$ has finite cohomological dimension. Then the
theorem applies without restriction on $\ast$, with $\cC$ the category
of complexes of injectives (respectively projectives).
\end{corollary}

\begin{definition}\label{VIII:4-2-5} % originally 2.4
Let $\cA$ and $\cB$ be two abelian categories and let $f:\cA\to \cB$ be
an additive functor.

We say that $f$ has \emph{finite cohomological dimension} on the right
(respectively left) if $\eR^+ f$ (respectively $\eL^- f$)
exists and there is an integer $m\geqslant 0$ such that, for every integer
$n>m$, $\eR^n f$ (respectively $\eL^{-n} f$) is zero on the objects
arising from $\cA$ under the functor $D$. The
\emph{cohomological dimension} on the right (respectively left) of $f$ is then the
least integer $m$ having this property. If $\eR^+ f$ (respectively
$\eL^- f$) exists and $f$ does not have finite cohomological dimension,
we say that $f$ has infinite cohomological dimension.

An object $X$ of $\cA$ is said to be \emph{$f$-acyclic} on the right
(respectively on the left) if, for every integer $n>0$,
$\eR^n f(D(X)) = 0$ (respectively $\eL^{-n} f(D(X))=0$).
\end{definition}

\begin{corollary}\label{VIII:4-2-3-2} % originally corollary 2
The data are those of Definition \ref{VIII:4-2-5}. Suppose there exists a
set $\sigma$ of objects of $\cA$, stable under direct sums, having the
following properties:
\begin{enumerate}
  \item Every acyclic complex
    $\in \operatorname{Ob}{\eK^+(\cA)}$ (respectively
    $\in\operatorname{Ob}{\eK^-(\cA)}$) of objects of $\sigma$ is taken by
    $f$ to an acyclic complex.
  \item Every object $X$ of $\cA$ embeds in (respectively is a quotient of)
    an object of $\sigma$. The hypotheses of Theorem \ref{VIII:4-2-3}
    hold with $\ast=+$ (respectively $\ast=-$) and with $\cC$ the category
    of complexes whose objects in every degree belong to $\sigma$. The
    elements of $\sigma$ are right (respectively left) $f$-acyclic
    objects.
\end{enumerate}

If, in addition, $f$ has finite right (respectively left) cohomological
dimension, the set of right (respectively left) $f$-acyclic objects has,
in addition to properties (1) and (2) above, property (2) of Proposition
\ref{VIII:3-1-9}. Theorem \ref{VIII:4-2-3} then applies to the functor
$\eK f$, with $\cC$ the category of complexes whose objects are right
(respectively left) $f$-acyclic in every degree. The functor $\eR f$
(respectively $\eL f$) exists and induces on $\eD^+(\cA)$ the functor
$\eR^+ f$ (respectively $\eL^+ f$), and on $\eD^-(\cA)$ the functor
$\eR^- f$ (respectively $\eL^- f$).
\end{corollary}

\subsection{Product. Composition}\label{VIII:4-3}

Let $\cA$ be an abelian category, let
$F:\eD(\cA) \to \mathsf{Ab}$ be a cohomological functor, and let
$X^\bullet$ and $Y^\bullet$ be two objects of $\eC(\cA)$. The functor $F$
defines a map
\[
  \ext^0(X^\bullet,Y^\bullet) \to \hom_{\mathsf{Ab}}(F(D(X^\bullet)),F(D(Y^\bullet)))
\]
and hence a pairing
$F(D(X^\bullet))\times \ext^0(X^\bullet,Y^\bullet) \to F(D(Y^\bullet))$.

Applying translations to this pairing gives pairings
\[
  F^i(D(X^\bullet))\times \ext^j(X^\bullet,Y^\bullet) \to F^{i+j}(D(Y^\bullet)) \text{.}
\]
By definition, when these pairings are applied to cohomological derived
functors, they are precisely the Yoneda products. The properties of these
products follow immediately from this definition. We shall not develop
them here.

Let $\cA$, $\cA'$, and $\cA''$ be three abelian categories, and let
$f:\cA\to \cA'$ and $g:\cA'\to \cA''$ be two additive functors. Suppose
that the functors $\eR^+ f$ and $\eR^+ g$ exist. I do not know how to
deduce that the functor $\eR^+ g\circ f$ exists, and even if the latter
functor exists, the formula
\[\xymatrix{
  \eR^+ g\circ f\ar[r]^-\sim
    & \eR^+ g \circ \eR^+ f
}\]
is probably not true in general. Nevertheless, one has the following
proposition.

\begin{proposition}\label{VIII:4-3-1}
Suppose that the category $\cA'$ has enough right $g$-acyclic objects and
that $\cA$ has enough right $f$-acyclic objects that $f$ takes to right
$g$-acyclic objects. Then the functor $\eR^+ g\circ f$ exists and there
is an isomorphism
\[\xymatrix{
  \eR^+ g\circ f \ar[r]^-\sim
    & \eR^+ g \circ \eR^+ f \text{.}
}\]
\end{proposition}

There are analogous statements for left derived functors.

\section{Examples}\label{VIII:6} % originally chapter 2, section 3

\subsection{The functor \texorpdfstring{$\mathsf{Rhom}^\bullet$}{Hom*}}\label{VIII:6-1}

Let $\cA$ be an abelian category with enough injectives.

The functor $Y\mapsto \hom^\bullet(X^\bullet,Y^\bullet)$
(\ref{VIII:1-3}), where $X^\bullet$ is an object of $\eK(\cA)$ and
$Y^\bullet$ an object of $\eK^+(\cA)$, has a total right derived functor
(Corollary \ref{VIII:4-2-3-1}), denoted by
$\rhom^\bullet(X^\bullet,-)$. Let $Y$ be an object of $\eD^+(\cA)$.
The functor
\begin{align*}
  \eK(\cA) &\to \eD(\mathsf{Ab}) \\
  X^\bullet &\mapsto \rhom^\bullet(X^\bullet,Y)
\end{align*}
is exact and vanishes on acyclic complexes, hence gives an exact
bifunctor
\[\xymatrix{
  \eD(\cA)^\circ \times \eD^+(\cA) \ar[r]
    & \eD(\mathsf{Ab})
}\]
which we denote by $\mathsf{Rhom}^\bullet$.

When $\cA$ has finite cohomological dimension, the functor
$\mathsf{Rhom}^\bullet$ extends to
$\eD(\cA)^\circ\times \eD(\cA)$ (loc. cit.).

If, in addition, $\cA$ has enough projectives, the functor $\hom$ can be
derived with respect to the first argument. The two definitions coincide
on their common domain of validity.

\subsection{\texorpdfstring{$\Tor$}{Tor} of sheaves}\label{VIII:6-2}

Let $E$ be a ringed site (in particular, a ringed topological space), let
$\sA$ be its sheaf of rings, and let $\mathsf{Mod}(E)$ be the category of
$\sA$-modules. An object $X$ of $\mathsf{Mod}(E)$ is said to be
\emph{$\sA$-flat} if, for every exact sequence
\[\xymatrix{
  0 \ar[r]
    & Y' \ar[r]
    & Y \ar[r]
    & Y'' \ar[r]
    & 0
}\]
the sequence
$0 \to Y'\otimes_\sA X \to Y\otimes_\sA X \to Y''\otimes_\sA X\to 0$
is exact. There are enough $\sA$-flat objects: the direct sums of objects
$\sA_U$, where $U$ is an object of $E$ and $\sA_U$ is the sheaf $\sA$
“extended by zero outside $U$.” Thus Theorem \ref{VIII:4-2-3} applies,
and one can define the functor
\begin{align*}
  \eD^-(E,\sA) \times \eD^-(E,\sA) &\to \eD^-(E,\sA) \\
  (X,Y) &\mapsto X\lotimes Y
\end{align*}
the total tensor product, which is the left derived functor of the tensor
product. Passing to cohomology defines the $\Tor_i^\sA(X,Y)$, the local
hyper-Tor groups.

Since the category $\mathsf{Mod}(E)$ has enough injectives, one can derive
the functor $\hom^\bullet(X,Y)$, the complex of local homomorphisms, and
thus obtain $\mathsf{Rhom}(X,Y)$ ($X\in \eD$, $Y\in \eD^+$).

Let $X$ and $Y$ be two objects of $\eD^-(E,\sA)$ and let $Z$ be an object
of $\eD^+(E,\sA)$. There is a functorial isomorphism
\[\xymatrix{
  \hom_{\eD(E,\sA)}(X\lotimes Y,Z) \ar[r]^-\sim
    &\hom_{\eD(E,\sA)}(X,\mathsf{Rhom}(Y,Z)) \text{.}
}\]

\subsection{Left derived functor of the inverse image}\label{VIII:6-3}

Let $f:E\to E'$ be a map of ringed topological spaces. The inverse-image
functor
\[
  f^\ast:\mathsf{Mod}(E') \to \mathsf{Mod}(E)
\]
is not necessarily exact. But one knows how to define $f^\ast$-flat
objects, and there are enough such objects. Thus one can define the
functor $\eL^- f^\ast$. There is likewise a duality formula:

Let $X$ be an object of $\eD^-(E',\sA')$ and let $Y$ be an object of
$\eD^+(E,\sA)$. There is a functorial isomorphism
\[\xymatrix{
  \hom_{\eD(E,\sA)}(\eL^- f^\ast(X),Y) \ar[r]^-\sim
    & \hom_{\eD(E',\sA')}(X,\eR^+ f_\ast(Y)) \text{.}
}\]
